\documentclass[11pt,a4paper,fontset=fandol]{ctexart}

\usepackage[a4paper,top=2.45cm,bottom=2.6cm,left=2.35cm,right=2.35cm,headheight=18pt,headsep=0.55cm,footskip=26pt]{geometry}
\usepackage{amsmath,amssymb,amsthm}
\usepackage{fancyhdr}
\usepackage{lastpage}
\usepackage{enumitem}
\usepackage{titlesec}
\usepackage{xcolor}
\usepackage{tikz}
\usepackage{hyperref}
\usepackage{bookmark}
\usepackage[most]{tcolorbox}
\usepackage{setspace}
\usepackage{array}
\usetikzlibrary{arrows.meta,positioning}

\hypersetup{
  colorlinks=true,
  linkcolor=blue!50!black,
  urlcolor=blue!50!black,
  citecolor=blue!50!black
}

\setstretch{1.17}
\setlength{\parindent}{2em}
\setlength{\parskip}{0.35em}
\allowdisplaybreaks
\setlist[enumerate]{itemsep=0.32em, topsep=0.32em, leftmargin=2.3em}
\setlist[itemize]{itemsep=0.25em, topsep=0.25em, leftmargin=2em}
\setcounter{tocdepth}{2}

\definecolor{mainblue}{RGB}{36,83,140}
\definecolor{lightblue}{RGB}{241,246,252}
\definecolor{lightgray}{RGB}{246,246,246}
\definecolor{accent}{RGB}{153,76,0}
\definecolor{rulegray}{RGB}{110,118,128}
\definecolor{softgreen}{RGB}{236,247,240}

\titleformat{\section}
  {\Large\bfseries\color{mainblue}}
  {\thesection.}
  {0.45em}
  {}
  [\vspace{0.25em}\color{rulegray}\titlerule]
\titlespacing*{\section}{0pt}{1.3em}{0.9em}
\titleformat{\subsection}{\large\bfseries\color{mainblue}}{\thesubsection}{0.5em}{}
\titlespacing*{\subsection}{0pt}{1.05em}{0.55em}
\titleformat{\subsubsection}{\normalsize\bfseries\color{accent}}{\thesubsubsection}{0.5em}{}

\pagestyle{fancy}
\fancyhf{}
\fancyhead[L]{\small 代数专题课堂讲义}
\fancyhead[C]{\small 函数图像对称与方程}
\fancyhead[R]{\small 刘文博}
\fancyfoot[C]{\small 第 \thepage 页 / 共 \pageref{LastPage} 页}
\renewcommand{\headrulewidth}{0.55pt}
\renewcommand{\footrulewidth}{0.35pt}

\newtheoremstyle{formaltheorem}
  {0.8em}{0.8em}{\normalfont}{}{\bfseries\color{mainblue}}{．}{0.6em}{}
\newtheoremstyle{formalexample}
  {0.8em}{0.8em}{\normalfont}{}{\bfseries\color{accent}}{．}{0.6em}{}

\theoremstyle{formaltheorem}
\newtheorem{theorem}{结论}[section]
\newtheorem{method}{方法}[section]
\theoremstyle{formalexample}
\newtheorem{example}{例题}[section]
\newtheorem{exercise}{练习}[section]

\newenvironment{solution}{\par\noindent\textbf{解：}}{\hfill$\square$\par}

\newtcolorbox{goalbox}[1]{
  breakable,
  colback=lightblue,
  colframe=mainblue,
  boxrule=0.7pt,
  arc=1mm,
  left=1.2mm,
  right=1.2mm,
  top=1mm,
  bottom=1mm,
  title=\textbf{#1},
  fonttitle=\bfseries
}

\newtcolorbox{notebox}[1]{
  breakable,
  colback=lightgray,
  colframe=black!50,
  boxrule=0.55pt,
  arc=1mm,
  left=1.2mm,
  right=1.2mm,
  top=1mm,
  bottom=1mm,
  title=\textbf{#1},
  fonttitle=\bfseries
}

\newtcolorbox{skillbox}[1]{
  breakable,
  colback=softgreen,
  colframe=green!40!black,
  boxrule=0.65pt,
  arc=1mm,
  left=1.2mm,
  right=1.2mm,
  top=1mm,
  bottom=1mm,
  title=\textbf{#1},
  fonttitle=\bfseries
}

\newcommand{\basef}{y=f(x)}

\begin{document}

\thispagestyle{empty}
\begin{titlepage}
\centering
\vspace*{0.9cm}

\begin{tcolorbox}[
  enhanced,
  width=0.92\textwidth,
  colback=white,
  colframe=mainblue,
  boxrule=1pt,
  arc=0mm,
  left=6mm,
  right=6mm,
  top=8mm,
  bottom=8mm
]
\centering

{\fontsize{24}{30}\selectfont\bfseries 代数专题课堂讲义\par}
\vspace{0.7cm}
{\fontsize{17}{22}\selectfont 函数图像对称与方程\par}

\vspace{1.1cm}
\color{rulegray}\rule{0.72\textwidth}{0.7pt}\par
\vspace{1cm}

\begin{tcolorbox}[colback=lightblue,colframe=mainblue,width=0.82\textwidth,boxrule=1pt]
\centering
{\Large 适合对象：初中阶段正在学习函数图像的学生}\\[0.4em]
{\large 已学过坐标系，接触过一次函数、二次函数或基本函数图像}
\end{tcolorbox}

\vspace{1.2cm}

\renewcommand{\arraystretch}{1.42}
\begin{tabular}{>{\bfseries}p{3.5cm}p{8cm}}
讲义作者： & 刘文博 \\
专题关键词： & 图像对称、关于 $x$ 轴、关于 $y$ 轴、关于原点、方程变形 \\
建议课时： & 1--2 课时 \\
专题目标： & 学会从已知函数方程快速写出对称后的新方程，并能反向判断对称方式 \\
编译方式： & XeLaTeX \\
\end{tabular}

\vspace{1.2cm}
\color{rulegray}\rule{0.72\textwidth}{0.7pt}\par
\vspace{0.8cm}

{\normalsize 本讲以“图像关于哪条轴或哪个点对称，方程应该怎么改”为主线，重点解决学生最容易混淆的“改 $x$ 还是改 $y$、在里面改还是在外面改”问题，并配上可直接订正的练习。\par}

\vspace{0.5cm}
{\large \today\par}
\end{tcolorbox}
\end{titlepage}

\pagenumbering{Roman}
\pagestyle{plain}
\tableofcontents
\newpage
\pagestyle{fancy}
\pagenumbering{arabic}

\section{专题目标与使用说明}

\begin{goalbox}{本讲要解决的核心问题}
已知一个函数的方程，例如
\[
y=f(x),
\]
如果把它的图像关于 $x$ 轴、关于 $y$ 轴、关于原点作对称，新的函数方程该怎么写？

这是函数图像学习里的关键一步，因为它把\textbf{图像变化}和\textbf{代数式变化}再次连在了一起，而且会直接影响以后学习二次函数、反比例函数和综合变换。
\end{goalbox}

\begin{goalbox}{学完以后希望达到的能力}
\begin{itemize}
  \item 看见“关于 $x$ 轴对称”就能立即写出 $y=-f(x)$；
  \item 看见“关于 $y$ 轴对称”就能立即写出 $y=f(-x)$；
  \item 看见“关于原点对称”就能立即写出 $y=-f(-x)$；
  \item 会用点坐标变化解释这些结论为什么成立；
  \item 会反向从新方程判断原图像经历了怎样的对称；
  \item 会处理平移和对称之间的简单衔接题。
\end{itemize}
\end{goalbox}

\begin{notebox}{给家长和自学同学的提醒}
\begin{itemize}
  \item 本讲的最大难点不在计算，而在\textbf{关于 $x$ 轴和关于 $y$ 轴容易混淆}；
  \item 建议每做一道题，都先口头说一遍：“横坐标变没变？纵坐标变没变？哪个变号了？”；
  \item 真正掌握的标准不是背公式，而是能用“图上的点怎样变”来解释每一个公式。
\end{itemize}
\end{notebox}

\section{先抓住最根本的点坐标变化}

\subsection{图像对称就是图上每个点一起变化}
如果点 $(x_0,y_0)$ 在图像 $\basef$ 上，那么一定有
\[
y_0=f(x_0).
\]
因此，研究图像对称时，最自然的方法就是研究：\textbf{图上的每一个点对称以后，新坐标变成什么。}

\begin{method}
若原图像上有点 $(x_0,y_0)$，那么：
\begin{itemize}
  \item 关于 $x$ 轴对称后，它变成 $(x_0,-y_0)$；
  \item 关于 $y$ 轴对称后，它变成 $(-x_0,y_0)$；
  \item 关于原点对称后，它变成 $(-x_0,-y_0)$。
\end{itemize}
\end{method}

\subsection{先分清谁变、谁不变}
\begin{itemize}
  \item 关于 $x$ 轴对称：横坐标不变，纵坐标变号；
  \item 关于 $y$ 轴对称：纵坐标不变，横坐标变号；
  \item 关于原点对称：横坐标和纵坐标都变号。
\end{itemize}

这三句话看起来简单，但几乎能决定整份讲义里所有结论。

\section{三种基本对称公式}

\subsection{关于 x 轴对称}

\begin{theorem}
已知原函数为
\[
y=f(x).
\]
把它的图像关于 $x$ 轴对称，所得新方程为
\[
y=-f(x).
\]
\end{theorem}

\begin{solution}
关于 $x$ 轴对称时，图像上每个点 $(x_0,y_0)$ 变成 $(x_0,-y_0)$。横坐标不变，只是函数值整体变成相反数，所以新方程就是
\[
y=-f(x).
\]
\end{solution}

\subsection{关于 y 轴对称}

\begin{theorem}
已知原函数为
\[
y=f(x).
\]
把它的图像关于 $y$ 轴对称，所得新方程为
\[
y=f(-x).
\]
\end{theorem}

\begin{solution}
关于 $y$ 轴对称时，图像上每个点 $(x_0,y_0)$ 变成 $(-x_0,y_0)$。纵坐标不变，横坐标变号。

设新图像上的横坐标记成 $x$，那么它对应原图像上的横坐标就是 $-x$，因此新图像上的函数值应为
\[
y=f(-x).
\]
\end{solution}

\subsection{关于原点对称}

\begin{theorem}
已知原函数为
\[
y=f(x).
\]
把它的图像关于原点对称，所得新方程为
\[
y=-f(-x).
\]
\end{theorem}

\begin{solution}
关于原点对称时，图像上每个点 $(x_0,y_0)$ 变成 $(-x_0,-y_0)$。这相当于：
\begin{itemize}
  \item 先关于 $y$ 轴对称，得到 $y=f(-x)$；
  \item 再关于 $x$ 轴对称，得到 $y=-f(-x)$。
\end{itemize}
所以关于原点对称后的方程为
\[
y=-f(-x).
\]
\end{solution}

\begin{skillbox}{一句话记忆}
\begin{itemize}
  \item 关于 $x$ 轴：\textbf{外面加负号}；
  \item 关于 $y$ 轴：\textbf{里面把 $x$ 变成 $-x$}；
  \item 关于原点：\textbf{里面外面都处理}，即 $y=-f(-x)$。
\end{itemize}
\end{skillbox}

\subsection{图像直观对比}

\begin{center}
\begin{minipage}{0.48\textwidth}
\centering
\begin{tikzpicture}[scale=0.72,>=Stealth]
  \draw[->,gray!70] (-3.2,0) -- (3.2,0) node[right] {$x$};
  \draw[->,gray!70] (0,-5.2) -- (0,5.2) node[above] {$y$};
  \draw[smooth,domain=-2.1:2.1,samples=100,thick,mainblue] plot (\x,{\x*\x});
  \draw[smooth,domain=-2.1:2.1,samples=100,thick,red!70!black] plot (\x,{-(\x*\x)});
  \node[mainblue] at (1.8,4.5) {$y=x^2$};
  \node[red!70!black] at (1.6,-3.8) {$y=-x^2$};
\end{tikzpicture}

\small 关于 $x$ 轴对称：外面加负号
\end{minipage}
\hfill
\begin{minipage}{0.48\textwidth}
\centering
\begin{tikzpicture}[scale=0.72,>=Stealth]
  \draw[->,gray!70] (-4.2,0) -- (4.2,0) node[right] {$x$};
  \draw[->,gray!70] (0,-1.2) -- (0,5.2) node[above] {$y$};
  \draw[smooth,domain=-3.2:2.2,samples=100,thick,mainblue] plot (\x,{\x+1});
  \draw[smooth,domain=-2.2:3.2,samples=100,thick,red!70!black] plot (\x,{-\x+1});
  \node[mainblue] at (2.2,3.5) {$y=x+1$};
  \node[red!70!black] at (-2.3,3.5) {$y=-x+1$};
\end{tikzpicture}

\small 关于 $y$ 轴对称：里面把 $x$ 变成 $-x$
\end{minipage}
\end{center}

\section{怎样快速判断一个函数有没有对称性}

\subsection{偶函数与关于 y 轴对称}

\begin{theorem}
如果一个函数满足
\[
f(-x)=f(x),
\]
那么它的图像关于 $y$ 轴对称。
\end{theorem}

\begin{solution}
若 $f(-x)=f(x)$，说明横坐标互为相反数的两个点，对应的函数值相同，也就是点 $(x,f(x))$ 和 $(-x,f(x))$ 同时在图像上，因此图像关于 $y$ 轴对称。
\end{solution}

\subsection{奇函数与关于原点对称}

\begin{theorem}
如果一个函数满足
\[
f(-x)=-f(x),
\]
那么它的图像关于原点对称。
\end{theorem}

\begin{solution}
若 $f(-x)=-f(x)$，说明点 $(x,f(x))$ 和 $(-x,-f(x))$ 同时在图像上，因此图像关于原点对称。
\end{solution}

\begin{notebox}{先有结论，再做验证}
初中阶段遇到“判断图像是否关于 $y$ 轴对称”或“是否关于原点对称”时，最常见的方法就是把 $x$ 换成 $-x$，看看函数值怎样变化。
\end{notebox}

\section{例题讲解}

\begin{example}
把抛物线
\[
y=x^2+3
\]
关于 $x$ 轴对称，求新方程。
\end{example}

\begin{solution}
关于 $x$ 轴对称，应在整个函数外面加负号，所以新方程为
\[
y=-(x^2+3)=-x^2-3.
\]
因此
\[
\boxed{y=-x^2-3}.
\]
\end{solution}

\begin{example}
把直线
\[
y=2x-1
\]
关于 $y$ 轴对称，求新方程。
\end{example}

\begin{solution}
关于 $y$ 轴对称，应把式子里的 $x$ 换成 $-x$，得到
\[
y=2(-x)-1=-2x-1.
\]
因此
\[
\boxed{y=-2x-1}.
\]
\end{solution}

\begin{example}
把函数
\[
y=x^3-2x
\]
关于原点对称，求新方程。
\end{example}

\begin{solution}
关于原点对称，应写成
\[
y=-f(-x).
\]
这里
\[
f(x)=x^3-2x.
\]
先求
\[
f(-x)=(-x)^3-2(-x)=-x^3+2x.
\]
再整体加负号：
\[
y=-(-x^3+2x)=x^3-2x.
\]
所以新方程仍然是
\[
\boxed{y=x^3-2x}.
\]
这说明它本身就关于原点对称。
\end{solution}

\begin{example}
已知函数
\[
y=(x+3)^2-4,
\]
说出它是由
\[
y=x^2
\]
经过哪些变化得到的，并写出它关于 $y$ 轴对称后的方程。
\end{example}

\begin{solution}
原式表示由 $y=x^2$ 先向左平移 $3$ 个单位，再向下平移 $4$ 个单位得到。

现在要关于 $y$ 轴对称，只需把式子里的 $x$ 换成 $-x$：
\[
y=((-x)+3)^2-4=(x-3)^2-4.
\]
所以对称后的方程是
\[
\boxed{y=(x-3)^2-4}.
\]
\end{solution}

\section{最容易犯的几个错误}

\begin{notebox}{误区一：把关于 $x$ 轴对称和关于 $y$ 轴对称混掉}
关于 $x$ 轴对称时，是纵坐标变号，所以写成
\[
y=-f(x).
\]
关于 $y$ 轴对称时，是横坐标变号，所以写成
\[
y=f(-x).
\]
一个是改外面，一个是改里面，千万别混。
\end{notebox}

\begin{notebox}{误区二：关于 $y$ 轴对称时，只改了一部分}
例如原函数是
\[
y=x^2-4x+1.
\]
关于 $y$ 轴对称后，不能只把前面的 $x^2$ 看成不变，而后面的 $-4x$ 随便改。

正确做法是把\textbf{整个式子里的 $x$ 都换成 $-x$}，即
\[
y=(-x)^2-4(-x)+1=x^2+4x+1.
\]
\end{notebox}

\begin{notebox}{误区三：关于原点对称时，漏掉一个负号}
关于原点对称是
\[
y=-f(-x),
\]
不是单独的 $f(-x)$，也不是单独的 $-f(x)$。必须同时处理里面和外面。
\end{notebox}

\section{课堂练习}

\begin{exercise}
把函数
\[
y=x^2
\]
关于 $x$ 轴对称，写出新方程。
\end{exercise}

\begin{exercise}
把函数
\[
y=|x|
\]
关于 $y$ 轴对称，写出新方程。
\end{exercise}

\begin{exercise}
把函数
\[
y=2x+3
\]
关于 $y$ 轴对称，写出新方程。
\end{exercise}

\begin{exercise}
把函数
\[
y=x^3
\]
关于原点对称，写出新方程。
\end{exercise}

\begin{exercise}
已知函数
\[
y=-x^2+4,
\]
说出它是由
\[
y=x^2-4
\]
经过怎样的对称得到的。
\end{exercise}

\begin{exercise}
已知函数
\[
y=(x-2)^2+1,
\]
写出它关于 $y$ 轴对称后的方程。
\end{exercise}

\begin{exercise}
判断下列函数图像是否关于 $y$ 轴对称，若是，请说明理由：
\[
y=x^4+2x^2+1.
\]
\end{exercise}

\begin{exercise}
判断下列函数图像是否关于原点对称，若是，请说明理由：
\[
y=x^5-3x.
\]
\end{exercise}

\section{练习解答}

\begin{solution}
\textbf{第 1 题：}
关于 $x$ 轴对称，就是在原函数外面加负号，所以
\[
\boxed{y=-x^2}.
\]
\end{solution}

\begin{solution}
\textbf{第 2 题：}
函数 $y=|x|$ 本来就关于 $y$ 轴对称，所以关于 $y$ 轴对称后仍然是
\[
\boxed{y=|x|}.
\]
\end{solution}

\begin{solution}
\textbf{第 3 题：}
关于 $y$ 轴对称，应把 $x$ 换成 $-x$，得到
\[
y=2(-x)+3=-2x+3.
\]
所以
\[
\boxed{y=-2x+3}.
\]
\end{solution}

\begin{solution}
\textbf{第 4 题：}
关于原点对称，应写成
\[
y=-f(-x).
\]
这里 $f(x)=x^3$，于是
\[
y=-(-x)^3=x^3.
\]
所以
\[
\boxed{y=x^3}.
\]
\end{solution}

\begin{solution}
\textbf{第 5 题：}
原函数是
\[
y=x^2-4.
\]
关于 $x$ 轴对称后，应变成
\[
y=-(x^2-4)=-x^2+4.
\]
因此题目中的函数是由原函数\textbf{关于 $x$ 轴对称}得到的。
\end{solution}

\begin{solution}
\textbf{第 6 题：}
关于 $y$ 轴对称，应把式子里的 $x$ 换成 $-x$：
\[
y=((-x)-2)^2+1=(x+2)^2+1.
\]
所以
\[
\boxed{y=(x+2)^2+1}.
\]
\end{solution}

\begin{solution}
\textbf{第 7 题：}
把 $x$ 换成 $-x$：
\[
f(-x)=(-x)^4+2(-x)^2+1=x^4+2x^2+1=f(x).
\]
因此这个函数满足
\[
f(-x)=f(x),
\]
所以图像\textbf{关于 $y$ 轴对称}。
\end{solution}

\begin{solution}
\textbf{第 8 题：}
把 $x$ 换成 $-x$：
\[
f(-x)=(-x)^5-3(-x)=-x^5+3x=-(x^5-3x)=-f(x).
\]
因此这个函数满足
\[
f(-x)=-f(x),
\]
所以图像\textbf{关于原点对称}。
\end{solution}

\section{本讲小结}

\begin{goalbox}{三个最关键的结论}
\[
\begin{array}{rcl}
\text{关于 }x\text{ 轴对称} &\Longrightarrow& y=-f(x),\\[0.4em]
\text{关于 }y\text{ 轴对称} &\Longrightarrow& y=f(-x),\\[0.4em]
\text{关于原点对称} &\Longrightarrow& y=-f(-x).
\end{array}
\]
\end{goalbox}

\begin{skillbox}{最后再提醒一次}
\begin{itemize}
  \item 关于 $x$ 轴：看外面；
  \item 关于 $y$ 轴：看里面；
  \item 关于原点：里面外面都要看；
  \item 遇到不会的题，就回到“图上的点怎样变”这个最根本的问题。
\end{itemize}
\end{skillbox}

\begin{notebox}{自检清单}
\begin{itemize}
  \item 我能不能不看答案，直接写出 $y=|x-2|+3$ 关于 $y$ 轴对称后的方程？
  \item 我能不能解释为什么“关于 $y$ 轴对称”要把 $x$ 换成 $-x$？
  \item 我会不会用 $f(-x)=f(x)$ 或 $f(-x)=-f(x)$ 判断图像对称性？
\end{itemize}
如果这三项都能做到，就说明这部分已经真正掌握了。
\end{notebox}

\end{document}
